\chapter{Vectors} \label{ch:vectors}

\begin{center}
    \begin{tabular}{ c|c }
    \textbf{Vector}                                    & \textbf{Scalar} \\
    \hline
    Has both \textbf{direction} and \textbf{magnitude} & Has only magnitude \\
    Displacement                                       & Length/distance \\
    Velocity                                           & Speed \\
    Gravitational fields                               & Mass
    \end{tabular}
\end{center}

\section{Notation}

\begin{subequations}
    2D:
    \begin{equation}
        \vec{v}
        = \bm{v}
        = \ab< v_x, v_y >
        = \ab( v_x, v_y )
        = v_x \bm{\hat{\imath}} + v_y \bm{\hat{\jmath}}
    \end{equation}
    3D:
    \begin{equation}
        \vec{v}
        = \bm{v}
        = \ab< v_x, v_y, v_z >
        = \ab( v_x, v_y, v_z )
        = v_x \bm{\hat{\imath}} + v_y \bm{\hat{\jmath}} + v_z \bm{\hat{k}}
    \end{equation}
\end{subequations}

Note that \(\vec{v}\) is more common in handwriting,
and \(\bm{v}\) is more common in print.

See \pgRef{sec:unit-vectors} for information about \(\bm{\hat{\imath}}\),
\(\bm{\hat{\jmath}}\), and \(\bm{\hat{k}}\).

\section{Magnitude}

\textbf{Magnitude} is the size of the vector.

\begin{subequations}
    2D:
    \begin{equation}
        \ab\| \bm{v} \| = \sqrt{{v_x}^2 + {v_y}^2}
    \end{equation}
    3D:
    \begin{equation}
        \ab\| \bm{v} \| = \sqrt{{v_x}^2 + {v_y}^2 + {v_z}^2}
    \end{equation}
\end{subequations}

\section{Unit Vectors} \label{sec:unit-vectors}

\textbf{Units vectors} are vectors with a magnitude of 1.
They are denoted using a "hat", e.g., \(\hat{v}\) is "v-hat".

The \textbf{standard unit vectors} are
\begin{equation}
    \begin{aligned}
        \bm{\hat{\imath}} &= \bm{\hat{x}} = \ab< 1, 0, 0 >, \\
        \bm{\hat{\jmath}} &= \bm{\hat{y}} = \ab< 0, 1, 0 >, \\
        \bm{\hat{k}} &= \bm{\hat{z}} = \ab< 0, 0, 1 >
    \end{aligned}
\end{equation}
\vfil

\section{Operations}

Vector addition:
\begin{equation}
    \bm{u} + \bm{v} = \ab< u_x + v_x, u_y + v_y, u_z + v_z >
\end{equation}

Scalar multiplication:
\begin{equation}
    c \bm{v} = \ab< c v_x, c v_y, c v_z >
\end{equation}

Dot product:
\begin{equation}
    \begin{aligned}
        \bm{u} \cdot \bm{v} &= u_x v_x + u_y v_y + u_z v_z \\
        &= \ab\| \bm{u} \| \ab\| \bm{v} \| \cos\theta
    \end{aligned}
\end{equation}

Cross product:
\begin{equation}
    \begin{aligned}
        \bm{u} \times \bm{v}
        &= \ab< u_y v_z - u_z v_y, u_z v_a - u_a v_z, u_a v_y - u_y v_a > \\
        &= \ab\| \bm{u} \| \ab\| \bm{v} \| \bm{\hat{w}} \sin\theta \\
        &= \bm{w}
    \end{aligned}
\end{equation}
where \(\bm{\hat{w}}\) is a unit vector orthogonal to both \(\bm{u}\) and \(\bm{v}\)
as given by the \textbf{right-hand rule}.
Hold your index, middle, and thumb fingers such that all three are perpendicular to each other.
Your index finger represents \(\bm{u}\), your middle finger \(\bm{v}\),
and your thumb \(\bm{w}\).
